%% shortname: Utility Maximization and Expenditure Minimization
\documentclass[11pt]{article}
\usepackage{microstyle}

\begin{document}

\lessonheader{7}{Chapter 3 (continued): continuity, utility maximization, and expenditure minimization}{}


\section{Continuity: no preference reversal}

Continuity has the implication that there is no preference reversal: the ranking
of two bundles cannot flip in the limit.

\topic{An equivalent definition}
$\succsim$ is continuous if and only if, for every bundle $y$, the superior and
inferior sets
\[ \{\,x\in X : x\succsim y\,\}\qquad\text{and}\qquad \{\,x\in X : y\succsim x\,\} \]
are \emph{closed}. Equivalently, their complements --- the sets with strict
preference, $\{\,x\in X : y\succ x\,\}$ and $\{\,x\in X : x\succ y\,\}$ --- are
\emph{open}. In words: there is no dramatic change in the ordering if $y$ changes
a bit.

\begin{figure}[ht]
\centering
\begin{tikzpicture}[scale=0.9,
    hint/.style={font=\footnotesize,align=left,inner sep=1pt},
    lead/.style={thin,black!55}]
  % ------------- panel (a): continuous preferences -------------
  \begin{scope}
    \fill[black!12] plot[domain=0.5:4.0,samples=60] (\x,{1.8/\x})
          -- (4.0,4.0) -- (0.5,4.0) -- cycle;
    \ecaxes{4.4}{4.4}
    \node[ecnode,anchor=south west] at (0,4.45) {$x_2$};
    \node[ecnode,anchor=west]      at (4.5,0)   {$x_1$};
    \draw[budget] plot[domain=0.5:4.0,samples=60] (\x,{1.8/\x});
    \node[bundle] at (1.34,1.34) {};
    \node[ecnode,anchor=north east] at (1.30,1.30) {$y$};
    \node[ecnode] at (2.6,3.3) {$\{x : x\succsim y\}$};
    \node[ecnode,anchor=south east] at (4.4,4.45) {$L=2$};
    % a sequence inside the set converging to a point on the boundary
    \foreach \d in {0.8,0.4,0.2,0.1} \node[bundle,inner sep=1.0pt] at ({2.4+\d},{0.75+\d}) {};
    \node[bundle] at (2.4,0.75) {};
    \node[ecnode,anchor=south east] at (3.15,1.6) {$z^{n}$};
    \node[ecnode,anchor=north east] at (2.38,0.72) {$z$};
    \node[hint,anchor=west] at (5.0,3.6) {\textbf{(a) Continuous preferences}};
    \draw[lead] (4.95,2.3) -- (3.28,1.58);
    \node[hint,anchor=west] at (5.0,2.3)
         {$z^{n}$ inside the shaded set: $z^{n}\succsim y$};
    \draw[lead] (4.95,1.0) -- (2.5,0.78);
    \node[hint,anchor=west] at (5.0,1.0)
         {limit $z$ on the curve: $z\sim y$,\\ so it is still in the set --- the\\
          boundary belongs to the set};
  \end{scope}
  % ------------- panel (b): lexicographic preferences ------------
  \begin{scope}[yshift=-6.0cm]
    \fill[black!12] (1.8,0) rectangle (4.0,4.0);
    \ecaxes{4.4}{4.4}
    \node[ecnode,anchor=south west] at (0,4.45) {$x_2$};
    \node[ecnode,anchor=west]      at (4.5,0)   {$x_1$};
    % ticks and guides for the worked example
    \draw[thin] (1.8,0.08) -- (1.8,-0.08) node[ecnode,below] {$1$};
    \draw[thin] (0.08,1.0) -- (-0.08,1.0) node[ecnode,left] {$1$};
    \draw[thin] (0.08,2.2) -- (-0.08,2.2) node[ecnode,left] {$2$};
    \draw[densely dotted,black!55] (0,1.0) -- (1.8,1.0);
    \draw[densely dotted,black!55] (0,2.2) -- (1.8,2.2);
    % the vertical line x_1 = 1: included above y, excluded below
    \draw[budget] (1.8,2.2) -- (1.8,4.0);
    \draw[dashed] (1.8,0) -- (1.8,2.2);
    \node[bundle] at (1.8,2.2) {};
    \node[ecnode,anchor=south east] at (1.75,2.25) {$y$};
    \draw[fill=white] (1.8,1.0) circle (0.09);
    \node[ecnode,anchor=south east] at (1.75,1.05) {$x$};
    \foreach \a in {3.2,2.5,2.1,1.95} \node[bundle,inner sep=1.0pt] at (\a,1.0) {};
    \draw[-{Stealth[length=3pt]},thin] (3.1,1.35) -- (2.1,1.35);
    \node[hint,align=center] at (2.9,2.55) {$x_1>1$:\\ $\succ y$};
    % callouts
    \node[hint,anchor=west] at (5.2,4.5) {\textbf{(b) Lexicographic preferences}};
    \draw[lead] (5.15,3.75) -- (1.88,3.7);
    \node[hint,anchor=west] at (5.2,3.75)
         {above $y$ on the line: tie in $x_1$,\\ more $x_2$ $\Rightarrow$ in the set};
    \draw[lead] (5.15,2.6) -- (3.28,1.08);
    \node[hint,anchor=west] at (5.2,2.6)
         {$x^{n}=(1+\tfrac1n,\,1)$: more $x_1$\\ $\Rightarrow x^{n}\succ y$ for every $n$};
    \draw[lead] (5.15,1.3) -- (4.3,0.85) -- (1.9,0.85);
    \node[hint,anchor=north west] at (5.2,1.6)
         {limit $x=(1,1)$ on the dashed part:\\ tie in $x_1$, less $x_2$ $\Rightarrow y\succ x$,\\
          so $x$ is \emph{not} in the set (open circle):\\ the ranking has reversed};
  \end{scope}
\end{tikzpicture}
\caption{\emph{Top:} with $L=2$ goods and continuous preferences the superior set
of $y$ (shaded) contains its boundary, the indifference curve through $y$: it is
closed, so a sequence $z^{n}$ of bundles at least as good as $y$ converges to a
bundle $z$ that is still at least as good as $y$. \emph{Bottom:} the worked example
below, with $y=(1,2)$ and $x^{n}=(1+\frac1n,1)\to x=(1,1)$ marked on the axes, for
lexicographic preferences (good 1 compared first). The
superior set of $y$ is everything to the right of the vertical line through $y$,
plus the part of that line \emph{above} $y$; the part below $y$ (dashed) is not in
it. The bundles $x^{n}$ lie to the right, so $x^{n}\succ y$ for every $n$, but their
limit $x$ lies on the dashed part, so $y\succ x$: a preference reversal. The
superior set is not closed.}
\end{figure}

\topic{Why lexicographic preferences reverse: a worked example}
Recall the sequence definition of continuity from Lesson 6: if $x^{n}\succsim y^{n}$
for every $n$, and $x^{n}\to x$, $y^{n}\to y$, then $x\succsim y$. A \emph{preference
reversal} is exactly a failure of this: the ranking holds all along the sequence
but flips at the limit.

Take lexicographic preferences with good 1 compared first,
\[ x\succ y \iff x_1>y_1,\ \text{or}\ x_1=y_1\ \text{and}\ x_2>y_2 . \]
Good 1 is infinitely more important than good 2: any extra amount of good 1, however
small, beats any amount of good 2. Good 2 only matters to break an exact tie in good
1. Now fix
\[ y=(1,\,2)\qquad(\text{so } y^{n}=y \text{ for all } n), \qquad
   x^{n}=\Bigl(1+\tfrac1n,\ 1\Bigr)\ \longrightarrow\ x=(1,\,1). \]
\begin{flatnum}
  \item \emph{Along the sequence.} For every $n$, $x^{n}_1=1+\frac1n>1=y_1$, so good 1
        decides and $x^{n}\succ y$. It does not matter that $x^{n}$ has less of good
        2 (1 against 2): the tiny surplus $\frac1n$ of good 1 outweighs it.
  \item \emph{At the limit.} $x=(1,1)$ has exactly as much good 1 as $y$, so good 1
        no longer decides. The tie is broken by good 2, and $1<2$, so $y\succ x$.
  \item \emph{The reversal.} $x^{n}\succ y$ for every $n$, yet $y\succ x$ for the
        limit. The surplus of good 1 shrinks smoothly to zero, but its
        \emph{effect} on the ranking does not shrink: it stays decisive right up to
        the limit and then vanishes all at once, handing the decision to good 2.
        Continuity says exactly that this cannot happen, so lexicographic
        preferences are not continuous.
\end{flatnum}

\runin{The same thing in terms of closed sets.} Every $x^{n}$ lies in the superior
set of $y$, but the limit $x$ does not. A set that does not contain
the limits of its own sequences is not closed, so the superior set of $y$ is not
closed --- in panel (b) of Figure 1, the missing piece is the dashed part of
the vertical line below $y$, where $x$ sits. With continuous preferences (panel
(a)) the boundary, the indifference curve, belongs to the set, and the limit of
bundles that are at least as good as $y$ is still at least as good as $y$.

\runin{Why the picture looks so odd.} Under lexicographic preferences no two
different bundles are indifferent, so there are no indifference curves at all: each
``indifference set'' is a single point. The same jump is behind the fact that no
utility function can represent them: a utility $u$ would have to give each vertical
line $\{x_1=a\}$ its own interval of values $\bigl(u(a,0),\,u(a,1)\bigr)$, with
the intervals of different lines not overlapping (everything on a line with more
good 1 ranks higher). That would be uncountably many disjoint intervals in $\R$,
each containing a different rational number --- impossible, since there are only
countably many rationals.

\section{Utility representation}

\topic{The main conclusion}
If $\succsim$ is rational and continuous, there exists a \emph{continuous} utility
function $u(x)$ that represents $\succsim$ (Debreu's theorem). There is no utility
function at all representing lexicographic preferences, which are rational but not
continuous.

\topic{Differentiability}
For calculus we need the extra assumption that $u(x)$ is differentiable. This is a
genuine restriction: it rules out, for example, Leontief (perfect-complement)
preferences,
\[ u(x_1,x_2)=\min\{x_1,x_2\}, \]
which are continuous but have a kink wherever $x_1=x_2$.

\topic{Other connections between $\succsim$ and $u(x)$}
\[ \succsim\ \text{is convex}\iff u(x)\ \text{is quasiconcave}, \]
that is, $u\bigl(\alpha y+(1-\alpha)z\bigr)\ge\min\{u(y),u(z)\}$ for all
$\alpha\in[0,1]$: every superior set $\{x:u(x)\ge c\}$ is convex. Likewise
$\succsim$ is strictly convex if and only if $u$ is strictly quasiconcave.


\section{Utility maximization}

Assume $u(x)$ exists and is continuous, and assume preferences are locally
non-satiated. The \emph{utility maximization problem} (UMP) is
\[ \max_{x\ge 0}\ u(x)\qquad\text{s.t.}\qquad p\cdot x\le w . \]

\begin{flatlist}
  \item \emph{Existence.} A solution exists, since $u(x)$ is continuous and, for
        $p\gg0$ and $w\ge0$, the budget set $B_{p,w}$ is non-empty and compact
        (Lesson 1), and a continuous function attains its maximum on a non-empty
        compact set (Weierstrass).
  \item The solution is the \emph{Marshallian} (or \emph{Walrasian}) demand
        $x(p,w)$.
  \item We may have many solutions, in which case we have a demand
        \emph{correspondence}.
\end{flatlist}

\topic{Some properties of $x(p,w)$}
\begin{flatlist}
  \item \emph{Walras' law}, $p\cdot x=w$ for every $x\in x(p,w)$: the budget
        binds, due to local non-satiation (Lesson 6).
  \item \emph{Homogeneity of degree zero} in $(p,w)$:
        $x(\alpha p,\alpha w)=x(p,w)$ for $\alpha>0$, because the feasible set
        does not change, $B_{\alpha p,\alpha w}=B_{p,w}$.
\end{flatlist}

\section{The indirect utility function}

The indirect utility function is the maximum value function of the UMP,
\[ v(p,w)=u\bigl(x(p,w)\bigr): \]
the highest utility one can achieve at prices $p$ and wealth $w$. Its properties:
\begin{flatlist}
  \item $v(p,w)$ is continuous in $(p,w)$ (for $p\gg0$, $w>0$), as a consequence
        of the more general \emph{theorem of the maximum}.
  \item $v(p,w)$ is homogeneous of degree zero in $(p,w)$, because the feasible set
        does not change.
  \item $v(p,w)$ is strictly increasing in $w$ (by local non-satiation, more
        wealth always allows a better bundle) and non-increasing in $p_i$, for all
        $i=1,\dots,L$ (a higher price only shrinks the budget set).
  \item $v(p,w)$ is \emph{quasi-convex} in $(p,w)$: the inferior set
        $\{(p,w) : v(p,w)\le\bar v\}$ is convex for every $\bar v$.
\end{flatlist}

\runin{Why quasi-convex.} Let $v(p,w)\le\bar v$ and $v(p',w')\le\bar v$, and put
$(p'',w'')=\alpha(p,w)+(1-\alpha)(p',w')$ with $\alpha\in[0,1]$. Any $x$ affordable
at $(p'',w'')$ satisfies
$\alpha(p\cdot x)+(1-\alpha)(p'\cdot x)\le\alpha w+(1-\alpha)w'$, so either
$p\cdot x\le w$ or $p'\cdot x\le w'$: $x$ was affordable in one of the two original
situations, hence $u(x)\le\bar v$. Taking the best such $x$,
$v(p'',w'')\le\bar v$.

\begin{figure}[ht]
\centering
\begin{tikzpicture}[scale=0.62]
  \fill[black!12] (0.05,0) -- plot[domain=0.05:4.2,samples=60] (\x,{1.6*sqrt(\x)})
        -- (4.2,0) -- cycle;
  \ecaxes{4.8}{4.2}
  \node[ecnode,anchor=south west] at (0,4.25) {$w$};
  \node[ecnode,anchor=west]      at (4.9,0)   {$p_i$};
  \draw[budget] plot[domain=0.05:4.2,samples=60] (\x,{1.6*sqrt(\x)});
  \node[ecnode,anchor=west] at (4.25,3.28) {$v(p,w)=\bar v$};
  \node[ecnode] at (2.9,1.1) {$\{v(p,w)\le\bar v\}$};
\end{tikzpicture}
\caption{Quasi-convexity of $v$, drawn in the $(p_i,w)$ plane with the other prices
fixed. Along the curve utility stays at $\bar v$; below it (less wealth, or a higher
price) utility is at most $\bar v$. That inferior set is convex, so the curve is
concave.}
\end{figure}

\topic{The intuition}
The slope of the curve $v(p,w)=\bar v$ is the extra wealth needed to keep utility
at $\bar v$ when $p_i$ rises: to first order, one unit of $p_i$ costs $x_i$, the
amount of good $i$ being bought. The lower the price was to start with, the more of
good $i$ is bought (along the curve utility is held at $\bar v$, so this is the
compensated response, which never goes the wrong way), so I need to be compensated more for a price change the lower
the price was to start with. The curve is steep at low $p_i$ and flattens out: it is
concave.


\section{More structure: uniqueness and differentiability}

Impose more structure by assuming that $u(x)$ --- which is \emph{exogenous}, a
primitive of the model --- is strictly quasiconcave. Then the solution is
\emph{unique}, and $x(p,w)$ is a demand \emph{function}. (Two different solutions
would give the same utility, and by strict quasiconcavity their average --- also
affordable, since $B_{p,w}$ is convex --- would give strictly more.)

Let us also assume that $x(p,w)$ and $v(p,w)$ are differentiable. Strictly speaking
this is ``a crime'': they are \emph{endogenous}, derived from the model, so their
properties should be proved rather than assumed. The good news is that we can now
use a wonderful and elegant theorem, the \emph{envelope theorem}.

\section{The envelope theorem and Roy's identity}

The envelope theorem tells us what happens to the maximum value, here the indirect
utility, when a parameter changes: the derivative of the value function equals the
partial derivative of the Lagrangian with respect to the parameter, evaluated at the
optimum. The UMP
\[ \max_{x}\ u(x)\qquad\text{s.t.}\qquad p\cdot x\le w \]
has the Lagrangian
\[ \mathcal{L}=u(x)-\lambda\,(p\cdot x-w). \]

\topic{The effect of wealth}
\[ \pd{v(p,w)}{w}
   =\pd{\mathcal{L}}{w}\bigg|_{\substack{x=x(p,w)\\ \lambda=\lambda^{*}}}
   =\lambda^{*}>0 . \]
The multiplier $\lambda^{*}$ is the \emph{marginal utility of wealth}; it is
positive because $v$ is strictly increasing in $w$.

\topic{The effect of a price}
\[ \pd{v(p,w)}{p_i}
   =\pd{\mathcal{L}}{p_i}\bigg|_{\substack{x=x(p,w)\\ \lambda=\lambda^{*}}}
   =-\lambda^{*}x_i(p,w)\le 0 . \]

\topic{Roy's identity}
Combining the two (dividing the second by the first) gives
\[ x_i(p,w)=-\,\frac{\partial v(p,w)/\partial p_i}{\partial v(p,w)/\partial w}, \]
so we can infer demand from the indirect utility function.


\section{The first-order conditions}

Start from the Lagrangian and assume the solution is interior, $x\gg0$. Then the
Karush--Kuhn--Tucker (KKT) conditions are
\[ \pd{u(x)}{x_i}-\lambda p_i=0
   \quad\iff\quad
   \pd{u(x)}{x_i}=\lambda p_i,\qquad i=1,\dots,L, \]
together with $\lambda\ge0$ and complementary slackness,
$\lambda\,(w-p\cdot x)=0$.

\topic{Is the constraint binding?}
If $\partial u(x)/\partial x_i>0$ for some good $i$ --- in particular if
$\nabla u(x)\gg0$ --- then $\lambda^{*}=\dfrac{\partial u/\partial x_i}{p_i}>0$, and
by complementary slackness $p\cdot x=w$: the constraint is binding, as Walras' law
says.

\topic{The tangency condition}
Dividing the conditions for $x_i$ and $x_j$, the multiplier cancels:
\[ \frac{\partial u(x)/\partial x_i}{\partial u(x)/\partial x_j}=\frac{p_i}{p_j}. \]
The marginal rate of substitution equals the relative price. This is the
\emph{tangency condition}: at the optimum the indifference curve is tangent to the
budget line.


\section{Expenditure minimization}

\topic{A thought experiment}
The \emph{expenditure minimization problem} (EMP) is
\[ \min_{x\ge0}\ p\cdot x\qquad\text{s.t.}\qquad u(x)\ge\bar u : \]
the cheapest way of getting utility no less than $\bar u$ (for a level
$\bar u>u(0)$). It is the \emph{dual} of utility maximization: we switched the
objective and the constraint. Instead of fixing the budget and reaching the highest
indifference curve, we fix the indifference curve ($\bar u$) and find the lowest
budget that reaches it.

\begin{figure}[ht]
\centering
\begin{tikzpicture}[scale=0.62]
  % ------------- panel (a): utility maximization -------------
  \begin{scope}
    \ecaxes{4.4}{4.4}
    \node[ecnode,anchor=south west] at (0,4.45) {$x_2$};
    \node[ecnode,anchor=west]      at (4.5,0)   {$x_1$};
    \draw[budget] (0,3.0) -- (3.6,0);
    \foreach \c in {1.3,2.7,4.4}
      \draw[thin] plot[domain={\c/4.2}:4.2,samples=60] (\x,{\c/\x});
    \node[bundle] at (1.8,1.5) {};
    \draw[-{Stealth[length=4pt]},thin] (0.55,0.55) -- (1.25,1.25);
  \end{scope}
  % ------------- panel (b): expenditure minimization ----------
  \begin{scope}[xshift=7.2cm]
    \ecaxes{4.4}{4.4}
    \node[ecnode,anchor=south west] at (0,4.45) {$x_2$};
    \node[ecnode,anchor=west]      at (4.5,0)   {$x_1$};
    \draw[budget] plot[domain={2.7/4.2}:4.2,samples=60] (\x,{2.7/\x});
    \node[ecnode,anchor=south] at (3.9,0.72) {$\bar u$};
    \draw[dashed] (0,3.9) -- (4.68,0);
    \draw[thin] (0,3.0) -- (3.6,0);
    \node[bundle] at (1.8,1.5) {};
    \draw[-{Stealth[length=4pt]},thin] (3.0,1.9) -- (2.55,1.45);
  \end{scope}
\end{tikzpicture}
\caption{The two problems reach the same tangency from opposite sides.
\emph{Left, utility maximization:} the budget line is fixed and the consumer moves out to the highest
indifference curve that touches it. \emph{Right, expenditure minimization:} the indifference curve $\bar u$ is
fixed and the budget line, keeping its slope $-p_1/p_2$, is shifted in until it just
touches the curve.}
\end{figure}

When the two problems are matched --- the EMP at $\bar u=v(p,w)$ --- they pick the
same bundle, and the minimum expenditure is exactly $w$.


\section{The expenditure function}

The expenditure function is the minimum value function of the EMP,
\[ e(p,\bar u)=p\cdot h(p,\bar u), \]
where $h(p,\bar u)$ denotes the solution of the EMP. Its properties (for $p\gg0$):
\begin{flatlist}
  \item $e(p,\bar u)$ is continuous in $(p,\bar u)$.
  \item $e(p,\bar u)$ is homogeneous of degree 1 in $p$: if all prices are scaled
        by $\alpha$, relative prices are unchanged, so the solution is unchanged,
        $h(\alpha p,\bar u)=h(p,\bar u)$. But the bundle now costs more:
        \[ e(\alpha p,\bar u)=\alpha\,e(p,\bar u),\qquad \alpha>0 . \]
  \item $e(p,\bar u)$ is increasing in $\bar u$ (a higher target costs more) and
        non-decreasing in $p_i$, for all $i=1,\dots,L$.
  \item $e(p,\bar u)$ is \emph{concave} in $p$.
\end{flatlist}

\begin{figure}[ht]
\centering
\begin{tikzpicture}[scale=0.62]
  \ecaxes{4.8}{4.4}
  \node[ecnode,anchor=south west] at (0,4.45) {Exp};
  \node[ecnode,anchor=west]      at (4.9,0)   {$p_i$};
  \draw[budget] plot[domain=0.12:4.4,samples=60] (\x,{1.9*sqrt(\x)});
  \node[ecnode,anchor=west] at (4.45,3.99) {$e(p,\bar u)$};
  \draw[dashed] (0.3,1.40) -- (4.4,4.58);
  \node[bundle] at (1.5,2.33) {};
  \draw[thin] (1.5,0.08) -- (1.5,-0.08);
  \node[ecnode,anchor=north] at (1.5,-0.08) {$\bar p_i$};
  \node[ecnode,anchor=south east] at (4.4,4.62) {$p\cdot h(\bar p,\bar u)$};
\end{tikzpicture}
\caption{The expenditure function is concave in each price. The dashed line is the
cost, at prices $p$, of the bundle $h(\bar p,\bar u)$ that is optimal at $\bar p$:
it is linear in $p_i$, touches $e$ at $\bar p$, and lies above it everywhere else,
because re-optimizing at the new prices can only lower the cost.}
\end{figure}

\topic{Proof of concavity}
Let $p''=\alpha p+(1-\alpha)p'$ with $\alpha\in[0,1]$, and let $x''$ denote the
solution of the expenditure minimization problem at prices $p''$. Then
\[
\begin{aligned}
  e(p'',\bar u)=p''\cdot x''
  &=\alpha\,p\cdot x''+(1-\alpha)\,p'\cdot x''\\
  &\ge\alpha\,e(p,\bar u)+(1-\alpha)\,e(p',\bar u).
\end{aligned}
\]
The inequality holds because $x''$ reaches utility $\bar u$, so it is a feasible
choice in the EMP at $p$ and at $p'$; the cheapest feasible bundle costs no more,
$p\cdot x''\ge e(p,\bar u)$ and $p'\cdot x''\ge e(p',\bar u)$. This confirms
concavity.

\begin{transcriptnotes}
  \item Page 1 of the original is a title page reading only ``Lesson 7''; it is
        not reproduced.
  \item The direct utility function is written $V(x)$ throughout the original,
        and the indirect utility function $V(P,w)$ on some pages and $v(P,w)$ on
        others. Both are typeset in the course (MWG) notation: $u(x)$ for utility
        and $v(p,w)$ for indirect utility. Prices are typeset in lower case $p$;
        the original writes $P$. The expenditure function, written $E(P,u)$, is
        typeset $e(p,\bar u)$, with the utility target written $\bar u$ so that it
        is not confused with the function $u(\cdot)$.
  \item ``equivilent'' $\to$ \emph{equivalent}; ``Walra's'' $\to$ \emph{Walras'};
        ``Homongenity'' $\to$ \emph{homogeneity}; ``continous'' $\to$
        \emph{continuous}; ``Lagrangian'', ``Kuhn--Tucker'', ``marginal'' and
        ``tangency'' are spelled out where the original abbreviates or misspells
        them (``managin'', ``tangent condition''); ``KKT'' is expanded as
        Karush--Kuhn--Tucker.
  \item \textbf{Corrected:} the set beside the first sketch on page 2 reads
        $\{x\in X \mid y\succ x\}$, which is the strict \emph{inferior} set of $y$.
        The shading is above the indifference curve, and the sentence it
        illustrates is about the superior set being closed, so the label is
        typeset as the superior set $\{x : x\succsim y\}$.
  \item The second sketch on page 2 is the lexicographic order with good 1 compared
        first (the vertical line through $y$ bounds the superior set), whereas the
        definition in Lesson 6 compares good 2 first. The labels
        ``$x^{n}\succ y^{n}$, $y\succ x$'' are read with $y^{n}=y$, as drawn; the
        figure shows the sequence $x^{n}\to x$. \textbf{Added:} a worked numerical example,
        $y=(1,2)$ and $x^{n}=(1+\frac1n,1)\to x=(1,1)$, showing step by step how the
        ranking $x^{n}\succ y$ flips to $y\succ x$ at the limit, tied to both the
        sequence and the closed-set definitions of continuity. Figure 1 is
        stacked vertically and annotated: the coordinates of the example on
        the axes, a converging sequence $z^{n}\to z$ in the continuous panel,
        and callouts explaining each region and point.
  \item ``Solution is manshallian on Walrasian demand'' $\to$ \emph{Marshallian
        or Walrasian demand}; ``demand corresponse'' $\to$ \emph{correspondence}.
  \item ``$V(P,w)$ is homogeity degree zero'' and ``$E(P,u)$ condition in
        $(P,u)$'' are read as \emph{homogeneous of degree zero} and
        \emph{continuous in $(p,\bar u)$}.
  \item The sentence beside the quasi-convexity sketch, ``I need to be compensated
        more for a price change the lower price was to start with'', is read as
        \emph{the lower the price was to start with}.
  \item ``Assume preferences are non-satiated'' carries an insertion mark with
        ``loc'' above it; it is read as \emph{locally non-satiated}.
  \item ``I'm pose more structure'' $\to$ \emph{Impose more structure}. The
        remark ``A crime'' after assuming $x(p,w)$ and $v(p,w)$ differentiable is
        kept, with the reason it is a crime (they are endogenous) spelled out.
  \item \textbf{Corrected:} the tangency condition divides
        $\partial V/\partial x_i$ by $\partial V/\partial x_i$ in the original;
        the denominator is $\partial u/\partial x_j$, which is what gives $p_j$ on
        the right.
  \item \textbf{Corrected:} the KKT condition is written
        $\partial V/\partial x_i=\partial P_i$ on the right of the ``$\iff$''; it
        is $\lambda p_i$.
  \item \textbf{Corrected:} ``$E(P,w)$ is concave in $P$'' is written with $w$ in
        place of the utility level; it is $e(p,\bar u)$.
  \item On page 10 the monotonicity property reads ``non-decreasing in $u$ and in
        $P_i$''. That is true, but in $\bar u$ the standard statement is stronger:
        with continuous, locally non-satiated preferences $e$ is \emph{strictly}
        increasing in $\bar u$; it is typeset as ``increasing''.
  \item \textbf{Added:} the definition of quasiconcavity and the strict-convexity
        counterpart; the reason the UMP has a solution (Weierstrass); a proof that
        $v$ is quasi-convex; the argument for uniqueness under strict
        quasiconcavity; the Lagrangian written out for the envelope theorem, with
        $\lambda^{*}$ named the marginal utility of wealth; complementary
        slackness; the reason for the inequality in the concavity proof (feasibility
        of $x''$); ``rational'' in the representation theorem and the name
        Debreu; the countable-rationals argument for why lexicographic preferences
        have no utility representation (``Why the picture looks so odd''); the
        kink as the reason Leontief is not differentiable; the first-order reading
        of the slope of $v(p,w)=\bar v$ as $x_i$ in ``The intuition''; the
        restrictions $p\gg0$, $w\ge0$ or $w>0$, and $\bar u>u(0)$; the remark that the UMP and EMP pick the same bundle when
        $\bar u=v(p,w)$; and, in the last figure, the dashed cost line of a fixed
        bundle, which shows why $e$ is concave.
  \item The sketches were redrawn: the superior set for $L=2$ and the lexicographic
        counter-example (page 2); the inferior set of $v$ in the $(p_i,w)$ plane
        (page 5); utility maximization against expenditure minimization (page 9);
        and $e(p,\bar u)$ against $p_i$ (page 10).
  \item Pages of the original, for tracing back: page 2 continuity, closed and
        open sets, and the two sketches; page 3 the representation theorem,
        lexicographic and Leontief preferences, convexity and quasiconcavity, and
        the UMP; page 4 existence, Marshallian demand, Walras' law, homogeneity and
        $v(p,w)$; page 5 the properties of $v$ and the quasi-convexity sketch;
        page 6 strict quasiconcavity, differentiability and the envelope theorem
        for $w$; page 7 the envelope theorem for $p_i$, Roy's identity and the KKT
        conditions; page 8 the binding constraint and the tangency condition;
        page 9 the EMP and its sketches; page 10 the expenditure function and its
        properties; page 11 the proof of concavity.
\end{transcriptnotes}

\end{document}
